\documentclass[11pt]{article}
\usepackage[margin=1in]{geometry}
\usepackage[T1]{fontenc}
\usepackage[utf8]{inputenc}
\usepackage{lmodern}
\usepackage{microtype}
\emergencystretch=2em
\usepackage{amsmath,amssymb}
\usepackage{booktabs}
\usepackage{graphicx}
\usepackage{xcolor}
\usepackage{etoolbox}
\usepackage[numbers,sort&compress]{natbib}
\usepackage[hidelinks]{hyperref}
\usepackage{cleveref}

% Auto-generated by scripts/make_tables.py -- do not edit.
\newcommand{\AuditOneNearIdentical}{0}
\newcommand{\AuditOneNearIdenticalTest}{0}
\newcommand{\AuditOnePhashConflicts}{119}
\newcommand{\AuditOnePhashPairs}{293}
\newcommand{\AuditOnePhashTrue}{0}
\newcommand{\AuditOneSimilar}{21}
\newcommand{\AuditOneSimilarConflicts}{5}
\newcommand{\AuditOneSimilarTest}{8}
\newcommand{\AuditOneUnique}{2280}
\newcommand{\AuditTwoNearIdentical}{5}
\newcommand{\AuditTwoNearIdenticalTest}{0}
\newcommand{\AuditTwoPhashConflicts}{61}
\newcommand{\AuditTwoPhashPairs}{176}
\newcommand{\AuditTwoPhashTrue}{5}
\newcommand{\AuditTwoSimilar}{25}
\newcommand{\AuditTwoSimilarConflicts}{2}
\newcommand{\AuditTwoSimilarTest}{9}
\newcommand{\AuditTwoUnique}{2279}
\newcommand{\CrossTaskShared}{1140}
\newcommand{\CrossTaskSplitMismatch}{1}
\newcommand{\NPrimaryTests}{10}
\newcommand{\PowerAccHalfWidth}{3.3}
\newcommand{\PowerAucHalfWidth}{0.036}
\newcommand{\PowerAucMde}{0.051}
\newcommand{\PowerClassHalfWidth}{4.6}
\newcommand{\PowerMcnemarMde}{5.1}
\newcommand{\ProbeOneBorderTest}{0.60}
\newcommand{\ProbeOneBorderVal}{0.68}
\newcommand{\ProbeOneGeometryTest}{0.50}
\newcommand{\ProbeOneGeometryVal}{0.50}
\newcommand{\ProbeOneIntensityTest}{0.78}
\newcommand{\ProbeOneIntensityVal}{0.76}
\newcommand{\ProbeOneThumbTest}{0.83}
\newcommand{\ProbeOneThumbVal}{0.82}
\newcommand{\ProbeTwoBorderTest}{0.66}
\newcommand{\ProbeTwoBorderVal}{0.68}
\newcommand{\ProbeTwoGeometryTest}{0.50}
\newcommand{\ProbeTwoGeometryVal}{0.50}
\newcommand{\ProbeTwoIntensityTest}{0.77}
\newcommand{\ProbeTwoIntensityVal}{0.71}
\newcommand{\ProbeTwoThumbTest}{0.76}
\newcommand{\ProbeTwoThumbVal}{0.74}
\newcommand{\ResultsAvailable}{no}


% Figures are produced by scripts/make_figures.py; absent ones render as a note.
\newcommand{\maybefigure}[3][0.95\linewidth]{%
  \IfFileExists{figures/#2.pdf}{\includegraphics[width=#1]{figures/#2.pdf}}{%
  \IfFileExists{figures/#2.png}{\includegraphics[width=#1]{figures/#2.png}}{%
  \fbox{\parbox{0.9\linewidth}{\centering\small\textit{Pending: #3 is generated by
  \texttt{scripts/make\_figures.py} once the corresponding runs exist.}}}}}}
\newcommand{\pendingnote}{\ifdefstring{\ResultsAvailable}{yes}{}{\par\noindent\textit{No
  test-set results have been generated yet; all result tables below are placeholders.}\par}}

\title{Vision--language models versus supervised image models for chest radiograph
classification on POLCOVID: a pre-registered, leakage-audited comparison of prompting,
retrieval-augmented in-context learning and fine-tuning}
\author{Author names withheld for review}
\date{}

\begin{document}
\maketitle

\begin{abstract}
\noindent\textbf{Background.} General-purpose and medical vision--language models (VLMs) can
classify medical images without task-specific training, but comparisons with supervised models
often rest on single runs, free-text answer parsing and test sets of unknown cleanliness.
\textbf{Methods.} On two balanced binary tasks derived from the multicentre POLCOVID chest
radiograph database (Task~1: normal vs.\ pneumonia; Task~2: non-COVID pneumonia vs.\ COVID-19;
$n=456$ test images each), we compare Gemma-3-4B and MedGemma-4B under zero-shot prompting,
in-context learning (ICL) with random or retrieved exemplars, and LoRA fine-tuning, against
convolutional networks trained from ImageNet or chest-radiograph initialisations and linear
probes on frozen encoders. VLMs are scored by the log-likelihood of each label rather than by
parsing generated text, which yields a continuous score for AUROC and allows label-bias
calibration. All comparisons, seeds, selection rules and multiplicity corrections were fixed
in an analysis plan before the test set was accessed; every test-set access is logged.
Before modelling, we audited the published splits at pixel level and probed for non-anatomical
shortcuts.
\textbf{Results.} The audit found no near-identical images across splits of Task~1 and
\AuditTwoNearIdentical{} near-identical pairs (none involving test images) in Task~2, but
\AuditOneSimilarTest{} and \AuditTwoSimilarTest{} test images have a highly similar training
image, and an $8\times8$-pixel thumbnail alone reaches a test AUROC of \ProbeOneThumbTest{}
(Task~1) and \ProbeTwoThumbTest{} (Task~2). Model results are reported in
\cref{sec:results} from automatically generated tables, together with the null and negative
findings.
\textbf{Conclusions.} The protocol, code, configurations, per-image predictions and test-access
log are released so that every number in this paper can be regenerated.
\end{abstract}

\section{Introduction}

Chest radiography is the most frequently performed imaging examination, and deep networks
trained on labelled radiographs reach high accuracy on in-distribution
benchmarks~\citep{rajpurkar2017chexnet,irvin2019chexpert}. A second route has appeared with
instruction-tuned vision--language models (VLMs), which can be prompted with an image and a
question, given a few labelled examples in context~\citep{brown2020language}, or adapted with
parameter-efficient fine-tuning~\citep{hu2022lora}. Medical variants such as
MedGemma~\citep{sellergren2025medgemma} are pretrained on medical images and text and are
proposed as generalist medical models~\citep{moor2023foundation}.

Whether these routes are competitive with supervised image models on a given task is an
empirical question, and the answer is easy to distort. COVID-19 radiograph classifiers have
repeatedly been shown to exploit dataset artefacts rather than pathology~\citep{degrave2021shortcuts,
roberts2021pitfalls}, and pneumonia classifiers to depend on the acquiring
site~\citep{zech2018variable}. VLM evaluations add their own sources of variance: ICL accuracy
depends on which exemplars are chosen and in which order~\citep{liu2022makes,lu2022fantastically},
models are biased towards particular labels~\citep{zhao2021calibrate,zhou2024batch}, and
string-matching generated answers discards the model's confidence and conflates formatting
failures with misclassifications.

This paper grew out of a set of exploratory notebooks that compared a ResNet-50, zero-shot and
few-shot Gemma-3-4B and a LoRA fine-tuned Gemma-3-4B on two POLCOVID-derived tasks
(\cref{sec:legacy}). Reviewing them revealed problems that make their numbers hard to interpret:
single runs, a prompt mismatch between fine-tuning and evaluation, free-text parsing, one
exemplar draw per shot count, and no audit of the splits. We rebuilt the study around a
pre-registered protocol. Our contributions are:
\begin{enumerate}
  \item A pixel-level audit of the published splits and a shortcut analysis showing how much
  class information survives in low-level, non-anatomical image statistics (\cref{sec:audit}).
  \item A like-for-like comparison of a general (Gemma-3) and a medical (MedGemma) VLM of the
  same architecture and size under zero-shot, ICL and LoRA adaptation, against five supervised
  baselines, with multiple seeds or exemplar draws per condition.
  \item Likelihood-based label scoring for VLMs, with contextual~\citep{zhao2021calibrate} and
  validation-estimated batch-style~\citep{zhou2024batch} correction of label bias, and an
  explicit comparison with the generation-and-parse protocol.
  \item Retrieval-augmented ICL that selects class-balanced exemplars by similarity in the
  embedding space of a radiograph-specific encoder (RAD-DINO,~\citealp{perezgarcia2025raddino}),
  with near-duplicates of the query excluded.
  \item A frozen statistical analysis plan with paired bootstrap inference, Holm correction and
  a power analysis that states which differences the fixed test set can resolve.
\end{enumerate}

\section{Related work}

\paragraph{Supervised chest radiograph classification.} CNNs such as
ResNet~\citep{he2016resnet} and DenseNet~\citep{huang2017densenet} trained on large labelled
datasets~\citep{irvin2019chexpert} are the standard approach; TorchXRayVision provides
DenseNets pretrained on several public radiograph datasets~\citep{cohen2022torchxrayvision}.
Self-supervised encoders trained only on radiographs, such as RAD-DINO, give strong frozen
features for linear probing~\citep{perezgarcia2025raddino}.

\paragraph{Shortcut learning and evaluation pitfalls.} \citet{degrave2021shortcuts} showed that
COVID-19 detectors trained on combined datasets rely on laterality markers, text and border
regions; \citet{roberts2021pitfalls} systematically reviewed COVID-19 imaging models and found none of
potential clinical use, citing problems such as duplicated images, unclear data provenance and
inadequate evaluation;
\citet{zech2018variable} showed that pneumonia models identify the source hospital. Reporting
guidelines CLAIM~\citep{mongan2020claim,tejani2024claim} and TRIPOD+AI~\citep{collins2024tripodai}
codify the information needed to judge such studies; our checklist is in the supplementary
material.

\paragraph{VLMs for medical images.} Gemma~3~\citep{gemma2025gemma3} is an open
instruction-tuned VLM family; MedGemma~\citep{sellergren2025medgemma} continues its
pretraining on medical data and is reported to improve chest radiograph classification over
Gemma~3 at the same size. Parameter-efficient fine-tuning with LoRA~\citep{hu2022lora}, often on
4-bit quantised weights~\citep{dettmers2023qlora}, makes adapting such models feasible on a
single GPU.

\paragraph{In-context learning and calibration.} The choice~\citep{liu2022makes} and
order~\citep{lu2022fantastically} of in-context exemplars changes accuracy substantially, and
retrieving exemplars similar to the query is a common remedy. Label bias can be corrected by
shifting scores by those obtained for a content-free input~\citep{zhao2021calibrate} or by the
mean score over a batch~\citep{zhou2024batch}. We transfer both ideas to image ICL, estimating the
batch statistic on validation rather than test data.

\section{Data}\label{sec:data}

\paragraph{Source.} POLCOVID~\citep{suwalska2023polcovid} contains 4,809 chest radiographs from
15 Polish hospitals acquired in 2020--2021, labelled normal, pneumonia (non-COVID) or COVID-19.
We use two derived Hugging Face datasets prepared for the original project,
\texttt{Henryayay/POLCOVID\_task1\_70\_10\_20} and \texttt{Henryayay/POLCOVID\_task2\_70\_10\_20},
pinned to fixed revisions (\texttt{src/cxr/tasks.py}). Each contains 2,280 images resized to
$512\times512$ RGB and split 70/10/20 into 1,596 training, 228 validation and 456 test images,
balanced exactly by class. Task~1 is normal vs.\ pneumonia, where the pneumonia class pools
COVID-19 and non-COVID pneumonia; Task~2 is non-COVID pneumonia vs.\ COVID-19. The positive
classes are pneumonia (Task~1) and COVID-19 (Task~2).

\paragraph{Limitations of the derivation.} The scripts that produced the derived datasets
(image selection, resizing, split assignment) are not available, and neither are patient or
hospital identifiers; the original POLCOVID metadata require registration. We therefore cannot
split by patient or hospital, cannot test external validity across sites, and cannot exclude
that images of the same patient fall into different splits. The pixel-level audit
(\cref{sec:audit}) bounds, but does not eliminate, this risk. The audit also shows that all
\CrossTaskShared{} pneumonia images of Task~1 reappear in Task~2 (half as non-COVID pneumonia,
half as COVID-19) with consistent split membership for all but \CrossTaskSplitMismatch{} image,
so the two tasks are not independent samples and results should not be pooled across them.

\section{Methods}\label{sec:methods}

\subsection{Dataset audit}\label{sec:audit-methods}

Every image was hashed (SHA-256) to find exact duplicates. Because perceptual hashes of chest
radiographs collide frequently between different patients (\cref{tab:audit}), near-duplicates
were found in two stages: all pairs within and across splits whose normalised $64\times64$
thumbnails had a Pearson correlation of at least 0.92 were taken as candidates and then
re-scored at $256\times256$. Pairs with $r\geq0.99$ are called \emph{near-identical} and pairs
with $r\geq0.95$ \emph{highly similar}. To quantify shortcut information, logistic regression
and gradient boosting classifiers (chosen on validation) were trained on non-anatomical features
only: image geometry, colour mode, global intensity statistics, statistics of a 12.5\% border
frame, and an $8\times8$ thumbnail. The audit ran before any model was evaluated
(\texttt{scripts/audit\_dataset.py}).

\subsection{Supervised baselines}

\paragraph{CNNs.} ResNet-50~\citep{he2016resnet} and DenseNet-121~\citep{huang2017densenet}
initialised from ImageNet, and the TorchXRayVision DenseNet-121 pretrained on public radiograph
datasets (\texttt{densenet121-res224-all},~\citealp{cohen2022torchxrayvision}), each with a new
binary head. Images were zero-padded to square (a no-op for these $512\times512$ images, kept
for generality) and resized to $224\times224$. Training used AdamW~\citep{loshchilov2019adamw}
(weight decay $10^{-4}$, batch size 32), one warm-up epoch followed by cosine decay, bfloat16
autocast, mild augmentation (affine $\pm7^\circ$, $\pm4\%$ translation, $\pm7\%$ scale;
brightness and contrast jitter $\pm10\%$), at most 30 epochs and early stopping after 7 epochs
without improvement in validation loss. The learning rate was chosen from
$\{10^{-5},10^{-4},3\cdot10^{-4}\}$ by validation loss using a separate sweep seed, and the
checkpoint with the lowest validation loss was kept. Each configuration was trained with five
seeds (0--4).

\paragraph{Linear probes.} L2-regularised logistic regression on frozen, L2-normalised
embeddings from RAD-DINO~\citep{perezgarcia2025raddino} and from an ImageNet ResNet-50, with the
regularisation strength chosen by validation negative log-likelihood. Probes are deterministic
given the embeddings, so a single seed is used.

\subsection{Vision--language models}

\paragraph{Models.} \texttt{gemma-3-4b-it}~\citep{gemma2025gemma3} and
\texttt{medgemma-4b-it}~\citep{sellergren2025medgemma}, which share architecture and tokenizer,
loaded in bfloat16 from pinned revisions.

\paragraph{Prompt.} The same instruction is used for every VLM condition, including the
fine-tuning targets: \emph{``You are an expert radiologist. Analyze this chest X-ray and classify
it as `$a$' or `$b$'. Output only the diagnosis.''}, where $a$ and $b$ are the class names. For
ICL, $k$ exemplars are presented as ``Example $i$:'' followed by the image and ``Diagnosis:
\emph{label}'', and then the query image follows with the instruction.

\paragraph{Likelihood scoring.} Instead of parsing generated text, we compute the
log-probability of each complete label continuation (label tokens followed by the end-of-turn
token) given the prompt, and use the positive-class logit
$s=\log p(\text{pos}\mid x)-\log p(\text{neg}\mid x)$. Both continuations share a cached prompt
prefix; on a small random Gemma-3 model the cached and uncached computations agree to about
$10^{-6}$ (\texttt{tests/test\_vlm\_scoring.py}). For the zero-shot and LoRA conditions we also
run greedy generation with the original string-matching rule to compare the two protocols on
the same images.

\paragraph{Zero-shot prompting.} The prompt above, without exemplars.

\paragraph{In-context learning.} $k\in\{2,4,6,8,10\}$ class-balanced exemplars from the
training split, in shuffled order. \emph{Random ICL} draws one exemplar set per draw and uses it
for all queries, with five independent draws. \emph{Retrieval ICL} chooses, for each query, the
$k/2$ training images of each class closest to the query in cosine similarity of RAD-DINO
embeddings, excluding candidates with similarity above 0.995 so that near-copies of the query
cannot be retrieved; five draws vary only the exemplar order. The value of $k$ reported for
each method is the one with the highest mean validation AUROC.

\paragraph{LoRA fine-tuning.} Adapters of rank 32 ($\alpha=32$, dropout 0.05) on language and
vision layers of the bfloat16 model (not 4-bit, unlike the preliminary runs), trained on the
training split with the instruction above and the label as target. Training ran for up to
10 epochs (learning rate $2\cdot10^{-4}$, effective batch 16) with evaluation after each epoch;
the epoch with the lowest validation loss was kept, with early stopping after 5 epochs without
improvement. Five seeds; the adapter was then merged and scored exactly like the zero-shot model.

\paragraph{Label-bias correction.} Three score variants are compared, all estimated without
test data: \emph{raw} $s$; \emph{contextual} $s-s_{\mathrm{cf}}$, where $s_{\mathrm{cf}}$ is the
score of the same prompt and exemplars with a uniform mid-grey query
image~\citep{zhao2021calibrate}; and \emph{validation-centred} $s-\bar{s}_{\mathrm{val}}$, where
$\bar{s}_{\mathrm{val}}$ is the mean score on the (balanced) validation split, a
validation-estimated form of batch calibration~\citep{zhou2024batch}. For a binary task with
balanced classes, both corrections shift the decision boundary but leave AUROC unchanged, so
they are evaluated with threshold metrics under the argmax rule.

\subsection{Further analyses}

\paragraph{Ensemble.} The best VLM and best supervised condition (by validation AUROC) are
combined by a weighted sum of validation-standardised scores, with the weight chosen on
validation.

\paragraph{Data efficiency.} ResNet-50, the linear probes and Gemma-3 LoRA are retrained on
class-stratified 10, 25 and 50\% subsets of the training split.

\paragraph{Robustness and shortcuts.} Every test image is also scored under six fixed
perturbations: blackening a 12.5\% border frame, blackening the central 70\% of the image
(which covers most of the lung fields), gamma 0.6 and 1.6, Gaussian noise ($\sigma=5\%$ of
range) and JPEG compression at quality 10. A model whose AUROC stays well above 0.5 with the
lungs masked must be using information outside them. For CNN checkpoints we compute
Grad-CAM~\citep{selvaraju2017gradcam} maps and report the fraction of attribution mass in the
central region (0.49 for a uniform map).

\subsection{Statistical analysis}\label{sec:stats}

The full plan is in \texttt{docs/analysis\_plan.md}; it was committed before the first test-set
evaluation, and \texttt{results/test\_access\_log.jsonl} records the commit of every evaluation
that touched the test split. Deviations are listed in the plan.

\paragraph{Endpoints.} The primary endpoint is test AUROC for the positive class. Secondary
endpoints are balanced accuracy, sensitivity, specificity, F1, MCC, AUPRC, Brier score and
expected calibration error (ECE, 10 bins; after validation temperature
scaling~\citep{guo2017calibration} in the main tables). Threshold metrics in the main tables use
the threshold maximising Youden's $J$~\citep{youden1950index} on validation.

\paragraph{Selection.} Hyperparameters, checkpoints, $k$ and ensemble weights are chosen on
validation only. When a comparison involves ``the best'' condition of a family, it is the one
with the highest mean validation AUROC.

\paragraph{Comparisons.} Five primary comparisons are made in each task (\NPrimaryTests{} tests):
best VLM vs.\ best supervised model; MedGemma vs.\ Gemma-3 after LoRA and zero-shot; retrieval
vs.\ random ICL for Gemma-3 (AUROC); and validation-centred vs.\ raw scores for Gemma-3 random
ICL (balanced accuracy under the argmax rule). The Holm procedure~\citep{holm1979simple} controls
the family-wise error rate at 0.05 over all \NPrimaryTests{}. Further comparisons are labelled
exploratory and are not adjusted.

\paragraph{Inference.} For conditions with several seeds or draws, the estimand is the
difference in the seed-averaged metric. Its 95\% CI and two-sided $p$-value come from a paired,
class-stratified bootstrap over test images~\citep{efron1993bootstrap} with 10,000 replicates;
each replicate resamples images once and recomputes every seed of both conditions on that
resample, so both image and seed variation are reflected. $p$-values use the add-one
correction. As supporting analyses we report DeLong's test~\citep{delong1988comparing} with the
fast algorithm of \citet{sun2014fast} on seed-averaged scores, and the exact McNemar
test~\citep{mcnemar1947note} on the first seed. Proportions have Wilson
intervals~\citep{wilson1927probable}.

\paragraph{Sample size.} The test set is fixed at 456 images. \Cref{tab:power} reports what it
can resolve: at an AUROC of 0.85, the 95\% CI half-width is about \PowerAucHalfWidth{}
(Hanley--McNeil,~\citealp{hanley1982meaning}), and the smallest AUROC difference detectable with
80\% power between models with score correlation 0.5 is about \PowerAucMde{}; for accuracy near
0.85, the CI half-width is about $\pm$\PowerAccHalfWidth{} points, and a paired McNemar
comparison~\citep{connor1987sample} with 15\% discordant pairs detects differences of about
\PowerMcnemarMde{} points. Differences below these values are reported as inconclusive, not as
evidence of equivalence.

\begin{table}[t]
\centering
\small
\caption{What the fixed test set ($n=456$; 228 per class) can resolve. CI: 95\% half-width of a single AUROC (Hanley--McNeil). MDE: minimum detectable difference between two correlated AUROCs at 80\% power, for inter-model score correlation $r$, at $\alpha=0.05$ and at the Holm worst case $\alpha/10$.}
\label{tab:power}
\begin{tabular}{ccccc}
\toprule
AUROC & CI half-width & MDE ($r{=}0.5$) & MDE ($r{=}0.5$, $\alpha/10$) & MDE ($r{=}0.7$)\\
\midrule
0.60 & 0.052 & 0.074 & 0.096 & 0.057 \\
0.70 & 0.048 & 0.068 & 0.089 & 0.053 \\
0.80 & 0.041 & 0.058 & 0.076 & 0.045 \\
0.85 & 0.036 & 0.051 & 0.066 & 0.039 \\
0.90 & 0.029 & 0.042 & 0.055 & 0.032 \\
0.95 & 0.021 & 0.030 & 0.039 & 0.023 \\
\bottomrule
\end{tabular}
\end{table}


\section{Results}\label{sec:results}
\pendingnote

\subsection{Dataset audit and shortcut information}\label{sec:audit}

\Cref{tab:audit} summarises the audit. No image appears twice within a task at the byte level
except one exact duplicate in Task~2 (\AuditTwoUnique{} unique of 2,280). Task~1 has no
near-identical pairs; Task~2 has \AuditTwoNearIdentical{}, none involving a test image. At the
looser threshold, \AuditOneSimilar{} (Task~1) and \AuditTwoSimilar{} (Task~2) pairs are highly
similar; \AuditOneSimilarTest{} and \AuditTwoSimilarTest{} test images respectively have such a
match in training, and \AuditOneSimilarConflicts{} and \AuditTwoSimilarConflicts{} pairs carry
different labels. Visual inspection (\cref{fig:audit}) suggests that some highly similar pairs
are follow-up images of the same patient, so a patient-level overlap between splits cannot be
ruled out; the leakage sensitivity analysis (\cref{tab:leakage}) removes these test images.
By contrast, a 64-bit perceptual hash at Hamming distance $\leq4$ flags \AuditOnePhashPairs{}
pairs in Task~1, \AuditOnePhashConflicts{} of them with different labels and
\AuditOnePhashTrue{} with $r\geq0.99$: hashing alone would have badly overstated duplication.

\begin{table}[t]
\centering
\scriptsize
\caption{Pixel-level audit of the published splits. Near-identical: Pearson $r\geq0.99$ between normalised $256\times256$ images; highly similar: $r\geq0.95$. ``Test'' columns count test images with a matching training image. The last column shows how many pairs a 64-bit perceptual hash flags at Hamming distance $\leq4$ (in parentheses: pairs with different labels), illustrating why hashing alone is unreliable for chest radiographs.}
\label{tab:audit}
\begin{tabular}{lrrrrrrrr}
\toprule
Task & Images & Unique & \multicolumn{2}{c}{Near-identical} & \multicolumn{3}{c}{Highly similar} & pHash$\leq$4 \\
 & & & pairs & test & pairs & label conflicts & test & pairs\\
\midrule
Task 1 (Normal vs.\ Pneumonia) & 2280 & 2280 & 0 & 0 & 21 & 5 & 8 & 293 (119) \\
Task 2 (Non-COVID pneumonia vs.\ COVID-19) & 2280 & 2279 & 5 & 0 & 25 & 2 & 9 & 176 (61) \\
\bottomrule
\end{tabular}
\end{table}


\begin{figure}[t]
\centering
\maybefigure[0.48\linewidth]{fig_audit_pairs_task1}{the Task 1 audit figure}\hfill
\maybefigure[0.48\linewidth]{fig_audit_pairs_task2}{the Task 2 audit figure}
\caption{Most similar cross-split image pairs found by the audit (left: Task~1, right: Task~2),
with Pearson correlation at $256\times256$ and labels.}
\label{fig:audit}
\end{figure}

Low-level image statistics carry substantial class information (\cref{tab:shortcut}). Image
geometry and colour mode carry none, because all images were resized to the same size and
format, but global intensity statistics reach a test AUROC of \ProbeOneIntensityTest{} and
\ProbeTwoIntensityTest{}, the border frame alone \ProbeOneBorderTest{} and \ProbeTwoBorderTest{},
and an $8\times8$ thumbnail \ProbeOneThumbTest{} and \ProbeTwoThumbTest{} in Tasks~1 and~2. Some
of this information is genuinely diagnostic (consolidation changes global intensity), but it
also reflects differences in acquisition between the hospitals that contributed each class,
which we cannot separate without site labels. These values are a floor that any model's AUROC
should be judged against.

\begin{table}[t]
\centering
\small
\caption{AUROC of simple classifiers trained on low-level, largely non-anatomical image features (validation-selected between logistic regression and gradient boosting). All images are $512\times512$ RGB, so geometry and colour mode carry no signal.}
\label{tab:shortcut}
\begin{tabular}{lcccc}
\toprule
Features & \multicolumn{2}{c}{Task 1} & \multicolumn{2}{c}{Task 2}\\
 & val & test & val & test\\
\midrule
Image geometry & 0.500 & 0.500 & 0.500 & 0.500 \\
Border statistics & 0.681 & 0.600 & 0.681 & 0.661 \\
Global intensity statistics & 0.760 & 0.783 & 0.709 & 0.768 \\
$8\times8$ thumbnail & 0.824 & 0.829 & 0.743 & 0.764 \\
Colour mode & 0.500 & 0.500 & 0.500 & 0.500 \\
\bottomrule
\end{tabular}
\end{table}


\subsection{Main comparison}

\Cref{tab:main_task1,tab:main_task2} report every condition at its validation-selected
configuration, as mean $\pm$ SD over seeds or exemplar draws; per-run values with DeLong
intervals are in \texttt{results/summary/run\_metrics.csv}. \Cref{tab:comparisons} gives the
pre-specified comparisons with paired bootstrap CIs, raw and Holm-adjusted $p$-values, and the
DeLong and McNemar supporting tests. \Cref{fig:roc} shows ROC and reliability curves.

\begin{table}[t]
\centering
\scriptsize
\caption{Main results, Task 1 (Normal vs.\ Pneumonia).}
\label{tab:main_task1}
\begin{tabular}{lccccccc}
\toprule
Method & $k$ & AUROC & Bal.\ acc. & Sens. & Spec. & ECE & runs\\
\midrule
\multicolumn{8}{l}{\textit{Pending: experiments not yet run; generated automatically once results exist.}}\\
\bottomrule
\end{tabular}
\par\vspace{2pt}\footnotesize Test-set mean\,$\pm$\,SD over seeds (supervised, LoRA) or exemplar draws (ICL). AUROC uses raw scores. Threshold metrics use the decision threshold maximising Youden's $J$ on validation; ECE is after validation temperature scaling. $k$: number of in-context examples, chosen on validation.
\end{table}

\begin{table}[t]
\centering
\scriptsize
\caption{Main results, Task 2 (Non-COVID pneumonia vs.\ COVID-19).}
\label{tab:main_task2}
\begin{tabular}{lccccccc}
\toprule
Method & $k$ & AUROC & Bal.\ acc. & Sens. & Spec. & ECE & runs\\
\midrule
\multicolumn{8}{l}{\textit{Pending: experiments not yet run; generated automatically once results exist.}}}}\\
\bottomrule
\end{tabular}
\par\vspace{2pt}\footnotesize Test-set mean\,$\pm$\,SD over seeds (supervised, LoRA) or exemplar draws (ICL). AUROC uses raw scores. Threshold metrics use the decision threshold maximising Youden's $J$ on validation; ECE is after validation temperature scaling. $k$: number of in-context examples, chosen on validation.
\end{table}

\begin{table}[t]
\centering
\scriptsize
\caption{Pre-registered comparisons on the test set. $\Delta$ is the difference in the seed-averaged metric (A$-$B) with a paired, class-stratified bootstrap CI over test cases (10{,}000 replicates). Conditions are selected on validation only. Primary $p$-values are Holm-adjusted across both tasks; secondary comparisons are exploratory.}
\label{tab:comparisons}
\begin{tabular}{llp{5.2cm}lccc}
\toprule
Tier & Task & Comparison (A vs.\ B) & Metric & $\Delta$ (95\% CI) & $p$ & $p_{\text{Holm}}$\\
\midrule
\multicolumn{7}{l}{\textit{Pending: final runs not yet run; generated automatically once results exist.}}\\
\bottomrule
\end{tabular}
\end{table}


\begin{figure}[t]
\centering
\maybefigure{fig_roc_reliability}{the ROC and reliability figure}
\caption{Test ROC curves (top) and reliability diagrams after validation temperature scaling
(bottom) for the validation-best condition of each method family, first seed or draw.}
\label{fig:roc}
\end{figure}

\subsection{In-context learning: exemplar selection, number of shots and label bias}

\Cref{tab:icl} and \cref{fig:icl} report ICL performance as a function of $k$ for random and
retrieved exemplars, with raw and corrected scores. The spread across draws measures sensitivity
to exemplar choice (random) or order (retrieval). The positive-prediction rate under the argmax
rule shows whether a condition collapses onto one label, as the preliminary Task~2 runs did
(\cref{tab:legacy}).

\begin{table}[t]
\centering
\scriptsize
\caption{Zero-shot and in-context learning on the test set as a function of the number of examples $k$ (mean\,$\pm$\,SD over 5 draws). Balanced accuracy uses the argmax decision ($\text{logit}\geq0$) on raw, validation-centred and contextually calibrated scores. ``Pos.\ rate'' is the fraction predicted positive; values near 0 or 1 indicate collapse onto one label.}
\label{tab:icl}
\begin{tabular}{clcccccc}
\toprule
Task & Method & $k$ & AUROC & Bal.\ acc.\ (raw) & Bal.\ acc.\ (val-centred) & Bal.\ acc.\ (contextual) & Pos.\ rate (raw)\\
\midrule
\multicolumn{8}{l}{\textit{Pending: ICL experiments not yet run; generated automatically once results exist.}}}}\\
\bottomrule
\end{tabular}
\end{table}


\begin{figure}[t]
\centering
\maybefigure{fig_icl_shots}{the ICL shots figure}
\caption{ICL test AUROC (left of each pair) and balanced accuracy under the argmax rule (right)
against the number of exemplars $k$; mean $\pm$ SD over five draws. Dashed lines: raw scores;
solid lines: validation-centred scores.}
\label{fig:icl}
\end{figure}

\subsection{Likelihood scoring versus generation}

\Cref{tab:generation} compares, on the same images, the original protocol (greedy generation and
string matching) with likelihood scoring, including the fraction of generations that matched
neither label.

\begin{table}[t]
\centering
\scriptsize
\caption{Original protocol (greedy generation + string matching) versus label-likelihood scoring with the argmax rule, on the same test images (mean over seeds).}
\label{tab:generation}
\begin{tabular}{clcccc}
\toprule
Task & Method & Parse failures & Acc.\ (generation) & Acc.\ (likelihood) & Agreement\\
\midrule
\multicolumn{6}{l}{\textit{Pending: final VLM runs not yet run; generated automatically once results exist.}}}}\\
\bottomrule
\end{tabular}
\end{table}


\subsection{Data efficiency}

\Cref{tab:efficiency} and \cref{fig:efficiency} show test AUROC when training on 10--100\% of
the training split.

\begin{table}[t]
\centering
\scriptsize
\caption{Test AUROC versus the fraction of the 1{,}596 training images used (class-stratified subsets; mean\,$\pm$\,SD over seeds).}
\label{tab:efficiency}
\begin{tabular}{clcccc}
\toprule
Task & Method & 10\% & 25\% & 50\% & 100\%\\
\midrule
\multicolumn{6}{l}{\textit{Pending: data-efficiency experiments not yet run; generated automatically once results exist.}}\\
\bottomrule
\end{tabular}
\end{table}


\begin{figure}[t]
\centering
\maybefigure[0.8\linewidth]{fig_efficiency}{the data-efficiency figure}
\caption{Test AUROC against the fraction of the training split used (class-stratified subsets).}
\label{fig:efficiency}
\end{figure}

\subsection{Robustness, shortcut reliance and leakage sensitivity}

\Cref{tab:robustness} and \cref{fig:robustness} report test AUROC under the six perturbations.
The centre mask is the most informative: AUROC well above 0.5 with the lung fields removed
indicates reliance on information outside them. \Cref{tab:leakage} repeats the main analysis
after removing test images with a highly similar training image.

\begin{table}[t]
\centering
\scriptsize
\caption{Test AUROC (mean over seeds) under test-time perturbations. ``centre\_mask'' removes the central 70\% of the image (most of the lung fields); AUROC well above 0.5 after this perturbation indicates reliance on non-pulmonary cues.}
\label{tab:robustness}
\begin{tabular}{clccccccc}
\toprule
Task & Method & Clean & border\_mask & centre\_mask & gamma\_low & gamma\_high & noise & jpeg\\
\midrule
\multicolumn{9}{l}{\textit{Pending: final runs not yet run; generated automatically once results exist.}}\\
\bottomrule
\end{tabular}
\end{table}

\begin{table}[t]
\centering
\scriptsize
\caption{Leakage sensitivity: test AUROC (mean over runs) with and without the test images that have a highly similar ($r\geq0.95$) training image.}
\label{tab:leakage}
\begin{tabular}{clcccc}
\toprule
Task & Method & $n$ excluded & AUROC (all) & AUROC (excl.) & max $|\Delta|$\\
\midrule
\multicolumn{6}{l}{\textit{Pending: final runs not yet run; generated automatically once results exist.}}\\
\bottomrule
\end{tabular}
\end{table}


\begin{figure}[t]
\centering
\maybefigure{fig_robustness}{the robustness figure}
\caption{Test AUROC under perturbation for the validation-best condition of each family.}
\label{fig:robustness}
\end{figure}

\section{Discussion}

\paragraph{What the design can and cannot show.} The protocol was built so that a reader can
trust the direction of each reported difference: selection is done on validation, test access
is logged, all conditions that were run are reported, and the analysis is fixed in advance.
The design cannot provide external validity. Both tasks come from one database whose classes
were collected at partly different hospitals, the derivation of the published splits is
undocumented, and patient identifiers are unavailable. The shortcut probes show that a
classifier on an $8\times8$ thumbnail already reaches an AUROC around
\ProbeOneThumbTest{}--\ProbeTwoThumbTest{}; a high AUROC on these tasks is therefore weak evidence
that a model has learned radiological signs of pneumonia or COVID-19. The robustness and
Grad-CAM analyses help to interpret which information models use but cannot settle the question
without site-stratified or external data.

\paragraph{Scoring and calibration of VLMs.} Likelihood scoring turns a VLM into a classifier
with a continuous score, which makes AUROC, calibration and threshold selection available and
separates discrimination from the label bias that dominated the preliminary ICL results. Both
corrections we study are a single additive shift of the logit: they change threshold metrics
under the argmax rule but cannot change AUROC, so any gain they bring is a gain in operating
point, not in discrimination.

\paragraph{Power.} With 456 test images, differences in AUROC smaller than about 0.05 cannot be
detected reliably (\cref{tab:power}), and fewer after correcting for \NPrimaryTests{}
comparisons. Several comparisons of interest, such as MedGemma vs.\ Gemma-3 after fine-tuning,
may fall in this range; we report them as inconclusive rather than as equivalence.

\section{Limitations}

\begin{itemize}
  \item \textbf{Single source, no patient or site identifiers.} No external test set, no
  patient-level split, and no adjustment for site. Highly similar cross-split pairs suggest that
  some patients contribute images to more than one split.
  \item \textbf{Undocumented derivation.} The derived datasets use 1,140 of the 2,426 normal
  images, and 1,140 of the 1,147 non-COVID pneumonia and of the 1,236 COVID-19 images in POLCOVID
  (570 of each for the Task~1 pneumonia class); how they were selected is unknown.
  \item \textbf{Overlapping tasks.} All pneumonia images of Task~1 are reused in Task~2, so the
  two tasks are not independent replications.
  \item \textbf{Model scale.} Only 4B-parameter VLMs were studied, for reasons of compute; larger
  models may behave differently.
  \item \textbf{Hyperparameters.} LoRA settings were not tuned (they follow the preliminary
  recipe), whereas the CNN learning rate was. This favours the CNNs slightly.
  \item \textbf{Labels.} Reference labels are those of the source database; we did not have
  access to reports or to RT-PCR results.
  \item \textbf{Clinical relevance.} Balanced binary tasks do not reflect clinical prevalence,
  and the study does not evaluate a clinical workflow. None of the models is suitable for
  clinical use.
\end{itemize}

\section{Ethics and data statement}

The study uses only publicly released, de-identified images from the POLCOVID
database~\citep{suwalska2023polcovid}; ethical approval for its collection and release is
described in that publication. No new data were collected
and no attempt was made to re-identify individuals. MedGemma is distributed under the Health AI
Developer Foundations terms of use, Gemma~3 under the Gemma terms of use, and both are used here
for research only. Code, configurations, per-image predictions, run metadata and the test-access
log are released with the paper; the images remain subject to the source database's terms.

\paragraph{Reproducibility.} \texttt{README.md} gives the commands to reproduce every table and
figure. Every run stores its configuration and hash, the git commit, the pinned dataset
revision and the software environment. Tables and figures are generated only by
\texttt{scripts/make\_tables.py} and \texttt{scripts/make\_figures.py} from \texttt{results/}.
CPU-only analyses (audit, probes, power, statistics) run on a laptop; VLM and CNN experiments
require a single GPU and are scripted in \texttt{notebooks/colab\_runner.ipynb}.

\paragraph{Reporting checklist.} The CLAIM~2024 and TRIPOD+AI items are mapped to sections of
this paper in \texttt{paper/checklist.md}.

\bibliographystyle{plainnat}
\bibliography{refs}

\appendix

\section{Preliminary results from the original notebooks}\label{sec:legacy}

The original notebooks (\texttt{legacy/}) each report one run per condition. \Cref{tab:legacy}
lists their accuracies, parsed automatically from the saved notebook outputs. We report them for
transparency, but they are not comparable with the results above, for the following reasons
found while reviewing the code:
\begin{itemize}
  \item The fine-tuning prompt omitted ``Output only the diagnosis.'', which the evaluation
  prompt contained.
  \item The few-shot notebooks passed already-templated text to the processor without disabling
  special tokens, so prompts began with two beginning-of-sequence tokens.
  \item Each shot count used a single exemplar draw, so the non-monotonic pattern in $k$ cannot
  be separated from exemplar-choice variance.
  \item Answers were obtained by generation and string matching; unmatched answers counted as
  errors.
  \item The VLM was loaded in 4-bit precision for both fine-tuning and evaluation.
  \item An early Task~2 notebook (v0) used the labels ``Pneumonia'' and ``COVID'' instead of the
  dataset's class names.
  \item The CNN checkpoint was selected by validation accuracy at a fixed learning rate, with
  random horizontal flips (which invert laterality) as augmentation; each condition was run
  once.
\end{itemize}
In Task~2, most ICL configurations predict almost exclusively one class (recall near~1 for one
class and near~0 for the other) and accuracy stays near 50\%, which motivated the label-bias corrections studied
here.

\begin{table}[t]
\centering
\scriptsize
\caption{Preliminary single-run results from the original notebooks, parsed automatically from their stored outputs (\texttt{scripts/extract\_legacy\_results.py}). These runs predate the protocol of this paper and have known issues (Section~\ref{sec:legacy}); they are reported for transparency, not as findings.}
\label{tab:legacy}
\begin{tabular}{clcccc}
\toprule
Task & Method & Notebook & Acc.\ (\%) & Recall$_{neg}$ & Recall$_{pos}$\\
\midrule
1 & ResNet-50 & orig. & 82.02 & 0.850 & 0.790 \\
1 & Gemma-3-4B LoRA & v1.0.0 & 87.94 & 0.900 & 0.860 \\
1 & Gemma-3-4B zero-shot & v1.0.2 & 57.02 & 0.290 & 0.850 \\
1 & Gemma-3-4B ICL, 2 shots & v1.0.0 & 54.61 & 0.140 & 0.950 \\
1 & Gemma-3-4B ICL, 2 shots & v1.0.2 & 56.36 & 0.190 & 0.930 \\
1 & Gemma-3-4B ICL, 4 shots & v1.0.2 & 64.69 & 0.540 & 0.750 \\
1 & Gemma-3-4B ICL, 6 shots & v1.0.2 & 70.83 & 0.710 & 0.710 \\
1 & Gemma-3-4B ICL, 8 shots & v1.0.2 & 64.04 & 0.560 & 0.720 \\
1 & Gemma-3-4B ICL, 10 shots & v1.0.2 & 64.47 & 0.740 & 0.550 \\
2 & ResNet-50 & orig. & 76.32 & 0.750 & 0.780 \\
2 & Gemma-3-4B LoRA & v1.0.0 & 70.83 & 0.780 & 0.640 \\
2 & Gemma-3-4B zero-shot & v1.0.0 & 50.22 & 1.000 & 0.000 \\
2 & Gemma-3-4B zero-shot & v1.0.2 & 50.00 & 1.000 & 0.000 \\
2 & Gemma-3-4B ICL, 2 shots & v1.0.0 & 50.00 & 1.000 & 0.000 \\
2 & Gemma-3-4B ICL, 4 shots & v1.0.0 & 49.78 & 0.920 & 0.080 \\
2 & Gemma-3-4B ICL, 6 shots & v1.0.0 & 50.00 & 1.000 & 0.000 \\
2 & Gemma-3-4B ICL, 8 shots & v1.0.0 & 50.44 & 0.540 & 0.470 \\
2 & Gemma-3-4B ICL, 10 shots & v1.0.0 & 49.56 & 0.970 & 0.020 \\
2 & Gemma-3-4B ICL, 2 shots & v1.0.2 & 54.61 & 0.190 & 0.900 \\
2 & Gemma-3-4B ICL, 4 shots & v1.0.2 & 52.85 & 0.930 & 0.130 \\
2 & Gemma-3-4B ICL, 6 shots & v1.0.2 & 50.00 & 0.990 & 0.010 \\
2 & Gemma-3-4B ICL, 8 shots & v1.0.2 & 50.00 & 1.000 & 0.000 \\
2 & Gemma-3-4B ICL, 10 shots & v1.0.2 & 50.22 & 1.000 & 0.000 \\
\bottomrule
\end{tabular}
\end{table}


\end{document}
